%\documentclass[11pt,a4paper,twoside]{article}

\usepackage[utf8]{inputenc}

\usepackage[T1,T2A]{fontenc}

\usepackage[english.ancient,russian.ancient]{babel}

\usepackage{graphicx}

\usepackage{array}

\usepackage[doi]{samgtu-bib}

\usepackage{samgtu}

\usepackage{samgtu-name}

\usepackage{mathtools}

\mathtoolsset{showonlyrefs}

\usepackage{desclist}

\usepackage{euscript}

\usepackage{mathrsfs} 

\usepackage{multicol}

\usepackage{mathrsfs}

\usepackage{cite}

\usepackage{cmap}

\usepackage{qrcode}

\allowdisplaybreaks[4]

\usepackage[nodayofweek]{datetime}

\usepackage[thinlines]{easybmat}



\renewcommand{\JournalYear}{2018}

\renewcommand{\JournalVolume}{22}

\renewcommand{\JournalNumber}{4}



\newdate{JournalDateSign}{29}{12}{2018}% Дата подписания Вестника в печать

\renewcommand{\JournalDateSign}{\displaydate{JournalDateSign}}

\newdate{JournalDatePrint}{16}{01}{2019}% Дата выхода Вестника из печати

\renewcommand{\JournalDatePrint}{\displaydate{JournalDatePrint}}

%\renewcommand{\JournalUPL}{16.56} % Усл. печ. л.

%\renewcommand{\JournalUIL}{16.53} % Уч.-изд. л.

\renewcommand{\JournalUPL}{15.24} % Усл. печ. л.

\renewcommand{\JournalUIL}{15.21} % Уч.-изд. л.

\renewcommand{\JournalRegNumber}{220/18} % Регистрационный номер

\renewcommand{\JournalOrderNumber}{2} % Номер заказа



%\renewcommand{\JournalRMonth}{Март} % Месяц выхода Вестника

%\renewcommand{\JournalEMonth}{March} % Месяц выхода Вестника

%\renewcommand{\JournalRMonth}{Июнь} % Месяц выхода Вестника

%\renewcommand{\JournalEMonth}{June} % Месяц выхода Вестника

%\renewcommand{\JournalRMonth}{Сентябрь} % Месяц выхода Вестника

%\renewcommand{\JournalEMonth}{September} % Месяц выхода Вестника

\renewcommand{\JournalRMonth}{Декабрь} % Месяц выхода Вестника

\renewcommand{\JournalEMonth}{December} % Месяц выхода Вестника



\newcommand{\totop}[1]{\raisebox{6pt}[1pt][2pt]{#1}}

\newcommand{\tottop}[1]{\raisebox{12pt}[1pt][2pt]{#1}} 

\newcommand{\totttop}[1]{\raisebox{18pt}[1pt][2pt]{#1}} 

\newcommand{\tottttop}[1]{\raisebox{24pt}[1pt][2pt]{#1}} 

\newcommand{\totttttop}[1]{\raisebox{30pt}[1pt][2pt]{#1}} 

\newcommand{\tobot}[1]{\raisebox{-6pt}[1pt][2pt]{#1}} 

\newcommand{\tobbot}[1]{\raisebox{-12pt}[1pt][2pt]{#1}} 

\newcommand{\tobbbot}[1]{\raisebox{-18pt}[1pt][2pt]{#1}}



\graphicspath{{./00PIC/}}

%

%\usepackage{samgtu-name}



\vsgtu{1627} %registration number



\FirstPage{380}

\LastPage{394}

%

%\pdfcrop

%\draft



\ResearchArticle



%\begin{document}



\hfill \qrcode[hyperlink,height=.8in]{http://doi.org/10.14498/vsgtu1627}

\vspace{-.8in}



\udc{517.956.4, 517.977.1}

\msc{80A23,  35K05,  93C20}



\rutitle{Модальная идентификация граничного воздействия~в~двумерной обратной задаче теплопроводности}

\rutitleshort{Модальная идентификация граничного воздействия\dots}



\entitle{Modal identification of a boundary input in~the~two-dimensional inverse heat conduction problem}

\entitleshort{Modal identification of a boundary input\dots}



\ruauthor{Э.~Я.~Рапопорт, А.~Н.~Дилигенская}{Р\,а\,п\,о\,п\,о\,р\,т~Э.~Я., Д\,и\,л\,и\,г\,е\,н\,с\,к\,а\,я~А.~Н.}

\enauthor{E.~Ya.~Rapoport, A.~N.~Diligenskaya}{R\,a\,p\,o\,p\,o\,r\,t~E.~Ya., D\,i\,l\,i\,g\,e\,n\,s\,k\,a\,y\,a~A.~N.}



\ruaffil{\rusamgtu.}



\enaffil{\ensamgtu.}



\ruauthorsdetails{2}{% Количество авторов

\fullauthor{\rupid{38448}{Эдгар Яковлевич Рапопорт}}

\orcid{0000-0002-0604-8801} % ORCID ID автора 

\regalia{доктор технических наук, профессор} 

\position{профессор} 

\dept{каф. автоматики и управления в~технических системах} 

\email{rapoport@samgtu.ru}

\fullauthor{\rupid{119447}{Анна Николаевна Дилигенская}}

\CAuthor % Автор-корреспондент

\orcid{0000-0002-9867-9781} % ORCID ID автора 

\regalia{кандидат технических наук, доцент} 

\position{доцент} 

\dept{каф. автоматики и управления в~технических системах} 

\email{adiligenskaya@mail.ru}



}



\enauthorsdetails{2}{

\fullauthor{\enpid{38448}{Edgar Ya. Rapoport}}

\orcid{0000-0002-0604-8801} % ORCID ID

\regalia{Dr. Tech. Sci.} 

\position{Professor} 

\dept{Dept. of Automatics and Control in Technical Systems}

\email[n]{rapoport@samgtu.ru}

\fullauthor{\enpid{119447}{Anna N. Diligenskaya}}

\CAuthor % Сorresponding Author

\orcid{0000-0002-9867-9781} % ORCID ID avtora 

\regalia{Cand. Tech. Sci.} 

\position{Associate Professor} 

\dept{Dept. of Automatics and Control in Technical Systems}

\email[n]{adiligenskaya@mail.ru}

}



\rukeywords{обратная задача теплопроводности, двумерная постановка, компактное множество, непрерывные и непрерывно дифференцируемые функции, модальная идентификация}

\enkeywords{two-dimensional inverse heat conduction problem, compact set, continuous and continuously differentiable functions, modal identification}



\ruabstract{Предлагается метод приближенного решения двумерной граничной обратной задачи теплопроводности на компактном множестве непрерывных вместе со своими первыми производными функций, позволяющий восстановить граничное воздействие, зависящее от времени и пространственной координаты.



Используется модальное описание объекта в форме бесконечной системы линейных дифференциальных уравнений относительно коэффициентов разложения температурного поля в ряд по собственным функциям исследуемой начально-краевой задачи. Такой подход приводит к восстановлению искомой величины плотности теплового потока в~виде взвешенной суммы конечного числа ее модальных составляющих. 

Их значения определяются по значениям временных мод температурного поля, которые находятся на основе его модального представления из экспериментальных данных. 

Использование математической модели объекта в~пространстве изображений по Лапласу и~метода конечных интегральных преобразований приводит к описанию идентифицируемых воздействий и температурного поля в форме их разложений в ряды по собственным функциям одинаковой пространственной размерности и формированию на этой основе замкнутой системы уравнений относительно искомых величин.



Решена задача планирования температурных измерений, обеспечивающая на линии контроля в конечный момент интервала идентификации минимизацию ошибки аппроксимации экспериментального температурного поля его модельным представлением в равномерной метрике оценивания температурных невязок. 



Предложенный подход позволяет построить последовательность приближений, равномерно сходящихся с увеличением числа учитываемых модальных составляющих, к искомому решению.



Численное решение задачи реализовано в среде имитационного моделирования динамических систем {\tt Simulink MATLAB}$^\circledR$  и показало удовлетворительную точность решения задачи. 



}



\enabstract{A method for the approximate solution of a two-dimensional inverse boundary heat conduction problem on a compact set of continuous and continuously differentiable functions is proposed. The method allows us to reconstruct a boundary action that depends on time and a spatial coordinate.



A modal description of the object is used in the form of an infinite system of linear differential equations with respect to the coefficients of the expansion of the state function in a series in eigenfunctions of the initial-boundary value problem under study. This approach leads to the restoration of the sought value of the heat flux density in the form of a weighted sum of a finite number of its modal components. Their values are determined from the temporal modes of the temperature field, which are found from the experimental data on the basis of the modal representation of the field.

 To obtain a modal description of the identified input and the temperature field in the form of their expansions into series in eigenfunctions of the same spatial dimension, the mathematical model of the object in the Laplace transform domain and the method of finite integral transformations are used. On this basis, a closed system of equations with respect to the unknown quantities is formed.

 

The proposed approach allows us to construct a sequence of approximations that uniformly converge to the desired solution with increasing number of considered modal components.



The problem of the temperature experimental design is solved. This solution ensures the minimization of the approximation error of the experimental temperature field by its model representation in the uniform metric of estimating temperature discrepancies on the control line at the final moment of the identification interval. 

}



\newdate{rapoporta}{14}{05}{2018}

\date{\displaydate{rapoporta}}

\newdate{rapoportb}{08}{06}{2018}

\revisiondate{\displaydate{rapoportb}}

\newdate{rapoportc}{11}{06}{2018}

\accepteddate{\displaydate{rapoportc}}



\newdate{rapoporte}{27}{06}{2018}

\renewcommand{\JournalDataOnline}{\displaydate{rapoporte}}





%\ForPeerReview

\makerutitle



%\newpage



\Section[n]{Введение}

Результаты решения обратных задач математической физики находят применение во многих областях техники, в том числе при решении задач идентификации и диагностики процессов 

теплообмена~\cite{ulg:bib:Ozisik, ulg:bib:Alifanov_Opening_Lecture, ulg:bib:Matsevityi_Param_funct_identif}. 

Теория обратных задач теплопроводности (ОЗТ)~\cite{ulg:bib:Ozisik, ulg:bib:Alifanov_Opening_Lecture, ulg:bib:Matsevityi_Param_funct_identif, ulg:bib:Alifanov_OZT, ulg:bib:Beck_Blackwel, ulg:bib:Vatulyan_IP_MDTT}, направленная на развитие вычислительных методов и подходов, позволяющих восстанавливать неизвестные и  недоступные для непосредственного измерения характеристики нестационарных тепловых процессов, в настоящее время является одним из базовых направлений современной технологической теплофизики.

Широкое применение находят граничные ОЗТ, позволяющие определять характеристики, входящие в граничные условия. 

Такие задачи, в том числе, рассматриваются в нелинейных постановках~\cite{ulg:bib:Alifanov_Russian_contribution, ulg:bib:Matsevityi_nonlinear_IHCP}, а также формулируются относительно многомерных областей~\cite{ulg:bib:Zhi_Qian_Numerical_solution, ulg:bib:Ching_yu_boundary_estimation, ulg:bib:Kuzin_2_D}.







Решение ОЗТ может основываться на одном из двух подходов, первый из которых состоит в аппроксимации оператора обратной задачи, не являющегося непрерывным, семейством непрерывных операторов и осуществляет регуляризацию некорректно поставленных задач~\cite{ulg:bib:Tikhonov_Arsenin}.

Второй подход реализует переход к условно-корректной постановке задачи на основе информации о~характере решения, 

что позволяет переопределить метрические пространства входных и выходных воздействий процесса. 

Такой подход предусматривает априорное назначение классов корректности, на которых задача становится устойчивой~\cite{ulg:bib:Ivanov, ulg:bib:Tikhonov}.



В данной работе предлагается основанный на результатах теории оптимального управления объектами с распределенными параметрами~\cite{ulg:bib:Butkovskiy} метод приближенного решения двумерной граничной ОЗТ на компактных множествах специальной структуры, не требующий применения численных регуляризирующих алгоритмов. 



\smallskip



\Section{Постановка двумерной граничной обратной задачи теплопроводности}

Рассматривается, подобно~\cite{ulg:bib:Kuzin_2_D, ulg:bib:Alifanov_OZT, ulg:bib:Reinhardt}, граничная обратная задача нестационарной теплопроводности в двумерной области, заданной декартовыми координатами $x$ и $y$, при одностороннем нагреве тела тепловым потоком $q(y,\varphi)$ и отсутствии теплообмена на ограничивающих плоскостях $x=0$, $y=0$ и $y=1$.

В этом случае базовая линейная математическая модель процесса описывается двумерным уравнением теплопроводности в~относительных единицах для изменяющегося во времени $\varphi $ температурного поля $\theta (x,y,\varphi )$:

\begin{equation}\label{eq:ulg:GeneralModel1}

\frac{\partial \theta (x,y,\varphi )}{\partial \varphi } = \frac{{\partial ^2}\theta (x,y,\varphi )}{\partial {x^2}} + \frac{{\partial ^2}\theta (x,y,\varphi )}{\partial {y^2}},

\quad 

 x, y \in ( 0, 1  ), \quad \varphi  \in ( 0,\varphi ^* )

\end{equation}

с начальным 

\begin{equation}\label{eq:ulg:BoundaryCondition2}

\theta (x,y,0) = 0,\quad x,y \in [0,\,1]

\end{equation}

и граничными условиями второго рода:

\begin{gather}\label{eq:ulg:BoundaryCondition2}

\frac{\partial \theta (x,y,\varphi )}{\partial x}\Bigr|_{x=0} = 

\frac{\partial \theta (x,y,\varphi )}{\partial y}\Bigr|_{y=0} = \frac{\partial \theta (x, y,\varphi )}{\partial y}\Bigr|_{y=1} = 0,\quad \varphi  \in [0,\,\varphi ^*],

\\

\label{eq:ulg:BoundaryCondition2}

\frac{\partial \theta (x,y,\varphi )}{\partial x}\Bigr|_{x=1} = q(y,\varphi ),\quad\varphi  \in [0,\,\varphi ^*].

\end{gather}



В граничной ОЗТ функция $q(y,\varphi)$ неизвестна и подлежит определению на основе дополнительной информации о температуре, измеряемой на некотором доступном для контроля подмножестве области $( x,y ) \in [ 0, 1]$ определения $\theta(x,y,\varphi)$.

При постановке двумерных ОЗТ в качестве такого подмножества обычно рассматривается некоторая линия 

$x = {x^*} = \rm const$, $0 \le {x^*} \le 1$, $y \in [ 0,1 ]$~\cite{ulg:bib:Kuzin_2_D, ulg:bib:Alifanov_OZT, ulg:bib:Reinhardt} или некоторое число точек, расположенных во всей области $x, y \in [ 0,1 ]$~\cite{ulg:bib:Ozisik, KolForKuz15}.



В настоящей работе в качестве входной информации для решения ОЗТ используются экспериментальные температурные зависимости $\theta ^*(x^*,{y_i},\varphi )$, полученные на интервале идентификации $\varphi  \in [ 0,\varphi ^* ]$  в~точках  $(x^*,y_i)$, $i = \overline {1,R}$, $x = x^* = \rm const$, $0 \le x^* \le 1$, по которым требуется восстановить неизвестное пространственно"=временное граничное воздействие  $q(y,\varphi)$ по температурной невязке, оцениваемой в равномерной метрике на интервале  $[ 0,\varphi ^* ]$.



\smallskip



\Section{Модальное представление объекта идентификации}

Для решения ОЗТ по восстановлению граничного управления $q(y,\varphi)$ будем использовать метод модальной идентификации, позволяющий представить искомое воздействие в виде разложения в усеченный ряд по собственным функциям краевой задачи~\cite{ulg:bib:Koshlyakov_Gliner_Smirnov}.

Модальное описание объекта с распределенными параметрами (ОРП)~\eqref{eq:ulg:GeneralModel1}--\eqref{eq:ulg:BoundaryCondition2} основывается на применении к модели объекта~\eqref{eq:ulg:GeneralModel1}--\eqref{eq:ulg:BoundaryCondition2} конечных интегральных преобразований~\cite{ulg:bib:Koshlyakov_Gliner_Smirnov}

\begin{equation}

\label{eq:ulg:IntegralPreobr}

\overline \theta  (\mu _m,\lambda _n,\varphi ) = \int _0^1 {\int _0^1 {\theta (x,y,\varphi )\,} } K(\mu _m,x) {K^*}(\lambda _n,y)\,dxdy 

\end{equation}

по пространственным переменным с~ядрами, равными собственным функциям  

$K(\mu _m,x) = \cos  \mu _m x $ и $K^*(\lambda _n,y) = \cos  \lambda _n y $  

объекта по соответствующему аргументу $x$  и $y$, где собственные числа  

$\mu _m^2 = {\pi ^2}{m^2}$, $m = 0, 1, 2, \dots$ 

и~$\lambda _n^2 = {\pi ^2}{n^2}$, $n = 0, 1, 2, \dots$  являются корнями уравнений 

$\sin \mu  = 0$, $\sin \lambda  = 0$  соответственно. 

Применение преобразования~\eqref{eq:ulg:IntegralPreobr} к модели ОРП~\eqref{eq:ulg:GeneralModel1}--\eqref{eq:ulg:BoundaryCondition2} позволяет перейти к описанию объекта в виде бесконечной системы линейных дифференциальных уравнений:

\begin{equation}\label{eq:ulg:System_LinearDifEq}

\frac{d\overline \theta  (\mu _m,\lambda _n,\varphi )}{d\varphi } =  - ( \mu _m^2 + \lambda _n^2 )\overline \theta  (\mu _m,\lambda _n,\varphi ) + {\overline q _n}(\varphi )\cos \mu _m ,\,\,m,n = 0, 1, 2, \dots

\end{equation}

относительно коэффициентов (временных мод) $\overline \theta  (\mu _m,\lambda _n,\varphi ) $ разложения функции $\theta (x,y,\varphi ) $ в двумерный сходящийся в среднем ряд по собственным функциям~\cite{ulg:bib:Koshlyakov_Gliner_Smirnov}:

\begin{gather}\label{eq:ulg:Infinite_Serial_Temperature}

\theta (x,y,\varphi ) = \sum\limits_{m = 0}^\infty  {\sum\limits_{n = 0}^\infty  {C_m^2C_n^2\overline \theta  (\mu _m,\lambda _n,\varphi )} } \cos \mu _m x \cos \lambda _n y,

\\

\bar \theta (\mu _n,\lambda _n,0) = {\bar \theta _0}(\mu _n,\lambda _n) = 0,\quad m, n = 0, 1, 2\dots\ .

\end{gather}

В~\eqref{eq:ulg:System_LinearDifEq} временные моды

\begin{equation}\label{eq:ulg:Mody_Control}

{\overline q _n}(\varphi ) = \overline q (\lambda _n,\varphi ) = \int _0^1 {q(y,\varphi )} \cos \lambda_n y\,dy,\quad n = 0, 1, 2\dots

\end{equation}

граничного управления  $q(y,\varphi)$ являются коэффициентами его разложения в~одномерный бесконечный ряд

\begin{equation}\label{eq:ulg:Infinite_Serial_Control}

q(y,\varphi ) = \sum\limits_{n = 0}^\infty  {C_n^2\overline q (\lambda _n,\varphi )} \cos {\lambda _n}y.

\end{equation}



Множители $C_m^2$  и  $C_n^2$ , нормирующие собственные функции, определяются следующим образом

\begin{gather}

\label{eq:ulg:NormirCoefficients_x}

C_m^2 = \biggl( \int _0^1 \cos ^2 \mu_mx \, dx  \biggr)^{ - 1} = \left\{ \begin{array}{cc}

1,&m = 0,\\

2,&m \ge 1;

\end{array} \right.

\\

\label{eq:ulg:NormirCoefficients_y}

C_n^2 = \biggl( \int _0^1 \cos ^2 \lambda_ny\,dy  \biggr)^{ - 1} = \left\{ \begin{array}{cc}

1,&n = 0,\\

2,&n \ge 1.

\end{array} \right.

\end{gather}



\Section{Метод модальной идентификации граничного управления}

Равенства~\eqref{eq:ulg:System_LinearDifEq} формально определяют искомые соотношения между модальными составляющими ${\overline q _n}(\varphi )$ и  $ \overline {\theta } (\mu _m,\lambda _n,\varphi ) $. 

Однако при учете любого достаточно большого конечного числа членов рядов~\eqref{eq:ulg:Infinite_Serial_Temperature} и~\eqref{eq:ulg:Infinite_Serial_Control} для $ m = \overline {0,M}$,

$n = \overline {0,N}$, $M$,~$N > 1 $, обеспечивающего требуемую точность описания температурного поля, соответствующая «укороченная» система $ (M + 1) \times (N + 1) $ уравнений~\eqref{eq:ulg:System_LinearDifEq} оказывается неразрешимой относительно меньшего числа  $ N + 1 $ неизвестных ${\overline q _n}(\varphi )$, которые должны быть найдены на основе поведения $ (M + 1) \times (N + 1) $  величин $ \overline {\theta } (\mu _m,\lambda _n,{\varphi }) $. 

Такая ситуация объясняется описанием $ \theta (x,y,\varphi)$ и $q(y,\varphi) $ рядами~\eqref{eq:ulg:Infinite_Serial_Temperature} и~\eqref{eq:ulg:Infinite_Serial_Control} различной пространственной размерности.



Для устранения этого затруднения будем, подобно~\cite{ulg:bib:Rapoport_Open-loop_controllability}, использовать модальное описание объекта на основе его модели~\eqref{eq:ulg:GeneralModel1}--\eqref{eq:ulg:BoundaryCondition2}, записанной в изображениях Лапласа по переменной  $ \varphi $:

\begin{gather}\label{eq:ulg:Model_Laplas}

p\widetilde \theta (x,y,p) = \frac{{\partial ^2}\widetilde \theta (x,y,p)}{\partial {x^2}} + \frac{{\partial ^2}\widetilde \theta (x,y,p)}{\partial {y^2}},\quad x,y \in ( 0,1 ),

\\

\label{eq:ulg:GU2_Laplas}

\frac{\partial \widetilde \theta (x,y,p)}{\partial x}\Bigr|_{x=1} = \widetilde q(y,p),

\\

\label{eq:ulg:GU2_Laplas}

\frac{\partial \widetilde \theta (x,y,p)}{\partial x}\Bigr|_{x=0} = 0,\quad 

\frac{\partial \widetilde \theta (x,y,p)}{\partial y}\Bigr|_{y=0} = 0,\quad 

\frac{\partial \widetilde \theta (x,y,p)}{\partial y}\Bigr|_{y=1} = 0,

\end{gather}

где $ p $ "--- комплексная переменная преобразования Лапласа; 

$\widetilde \theta (x,y,p)$, $\widetilde q(y,p)$ "--- изображения функций  $ \theta (x,y,\varphi)$ и  $ q(y,\varphi)$.



Применение конечного интегрального преобразования  по пространственному аргументу $ y$ граничного управляющего воздействия  $ \widetilde q(y,p)$ к уравнениям объекта~\eqref{eq:ulg:Model_Laplas}--\eqref{eq:ulg:GU2_Laplas} приводит к обыкновенным дифференциальным уравнениям 

\begin{equation}\label{eq:ulg:OrdDifEq_Laplas}

\frac{{d^2} {\widetilde {\overline \theta  }}_n}{d{x^2}} - A_n^2{\widetilde {\overline \theta  }} = 0,\quad 

A_n^2 = p + \lambda _n^2,\quad n = 0, 1, 2, \dots

\end{equation}

по переменной  $ x $ относительно модальных составляющих 

\begin{equation}\label{eq:ulg:TemperMody_Laplas}

{\widetilde {\overline {\theta } }_n}(x,p) = \int _0^1 {\widetilde {\theta }(x,y,p)}  \cos \lambda_n y\,dy,\quad n = 0, 1, 2,\dots

\end{equation}

разложения

\begin{equation}\label{eq:ulg:Series_Temper_Laplas}

\widetilde \theta (x,y,p) = \sum\limits_{n = 0}^\infty  {C_n^2{{\widetilde  {\overline {\theta }} }_n}(x,p)}  \cos \lambda_ny

\end{equation}

 в ряд по ортогональной системе собственных функций $ \cos \lambda_ny $, дополняемым граничными условиями

\begin{equation}

\label{eq:ulg:GU_TemperMody_Laplas}

\frac{d{{\widetilde  {\overline \theta } }_n}(x,p)}{dx}\Bigr|_{x=0} = 0,\quad \frac{d{{\widetilde  {\overline \theta } }_n}(x,p)}{dx}\Bigr|_{x=1} = {\widetilde  {\overline q} _n}(p).

\end{equation}

Здесь $ {\widetilde  {\overline q} _n}(p) $ "--- модальные составляющие граничного воздействия, равные коэффициентам

\begin{equation}\label{eq:ulg:Coeff_GU_Laplas}

{\widetilde  {\overline q} _n}(p) = \int _0^1 {\widetilde q(y,p)}  \cos \lambda_n y\,dy,\quad n = 0, 1, 2, \dots

\end{equation}

 разложения 

\begin{equation}\label{eq:ulg:Series_GU_Laplas}

\widetilde q(y,p) = \sum\limits_{n = 0}^\infty  {C_n^2{{\widetilde  {\overline q} }_n}(p)}   \cos \lambda_n y

\end{equation}

в ряд по собственным функциям  $  \cos \lambda_n y $.



Решение задачи~\eqref{eq:ulg:OrdDifEq_Laplas},~\eqref{eq:ulg:GU_TemperMody_Laplas} на линии  $  x=x^* $ температурных измерений находится известными методами в виде

\begin{equation}

%\label{eq:ulg:TemperMody_via_GU_Mody_Laplas}

\widetilde  {\overline \theta } _n(x^*,p) =  H_n(x^*,p){\widetilde  {\overline q} _n}(p), \quad 

%\end{equation}

% \begin{equation}

 \label{eq:ulg:TransferFunc_TemperMody_GU_Mody_Laplas}

 H_n(x,p) = \frac{\ch\bigl({A_n}(p) \cdot x\bigr)}{{A_n}(p) \sh\bigl({A_n}(p)\bigr)},

\end{equation}

однозначным образом связывающем в операторной форме моды  ${\widetilde  {\overline q} _n}(p) $ \linebreak и ${\widetilde  {\overline \theta } _n}(x^*,p) $   разложений $  \widetilde q(y,p)$  и  $\widetilde \theta (x^*,y,p) $ в одномерные ряды~\eqref{eq:ulg:Series_Temper_Laplas},~\eqref{eq:ulg:Series_GU_Laplas} по пространственной координате $ y $   при учете любого числа $  N $  составляющих в~\eqref{eq:ulg:Series_Temper_Laplas} и~\eqref{eq:ulg:Series_GU_Laplas}

\begin{equation}\label{eq:ulg:GU_Mody_via_TemperMody_Laplas}

{\widetilde  {\overline q} _n}(p) =  H_n^{ - 1}(x^*,p){\widetilde  {\overline \theta } _n}(x^*,p),\quad n = \overline {0,N} .    

\end{equation}



Таким образом, предлагаемый метод решения условно-корректной задачи идентификации граничного управления, рассматриваемого на компактном множестве непрерывных вместе со своими первыми производными модальных переменных   $  {\overline \theta  _n}(x^*,\varphi )$, реализуется в пространстве изображений по Лапласу и состоит в восстановлении искомой величины $  \widetilde q(y,p)$  в виде конечной взвешенной суммы

\begin{equation}\label{eq:ulg:Finit_Series_GU_Laplas}

\widetilde q(y,p) = \sum\limits_{n = 0}^N {C_n^2{\widetilde  {\overline q} _n}(p)}   \cos \lambda_n y    

\end{equation}

мод ${\widetilde  {\overline q} _n}(p) $, вычисляемых согласно~\eqref{eq:ulg:TransferFunc_TemperMody_GU_Mody_Laplas},~\eqref{eq:ulg:GU_Mody_via_TemperMody_Laplas}, по значениям модальных составляющих  $ {\widetilde  {\overline \theta } _n}(x^*,p) $ температурного поля  $  \widetilde \theta (x^*,y,p) $.  

Выражения  $ {\widetilde  {\overline \theta } _n}(x^*,p)   $  могут рассматриваться как преобразованные по Лапласу зависимости  $  {\overline \theta  _n}(x^*,\varphi ) $, значения которых должны быть найдены по экспериментальным данным  путем решения относительно  $  {\overline \theta  _n}(x^*,\varphi ) $  замкнутой системы уравнений, образуемой оригиналами равенств~\eqref{eq:ulg:Series_Temper_Laplas}, рассматриваемых в точках контроля 

$x = {x^*}$, $y = {y_i}$, $i = \overline {1,N} ,$ для укороченной суммы ряда

\begin{equation}\label{eq:ulg:Finit_Series_Experiment}

{\theta ^*}(x^*,{y_i},\varphi ) = \sum\limits_{n = 0}^N {C_n^2{{\overline \theta  }_n}(x^*,\varphi )}  \cos \lambda_n{y_i},\quad i = \overline {1,N + 1}.    

\end{equation}



Предварительно на основе преобразования Лапласа должен быть выполнен переход от температурных кривых  $  {\theta ^*}(x^*,y_i,\varphi )  $, зависящих от времени, к их изображениям  $ \widetilde \theta ^*(x^*,y_i,p) $. Переход к оригиналу $ q(y,\varphi )  $  производится применением обратного преобразования Лапласа к равенству~\eqref{eq:ulg:Finit_Series_GU_Laplas} с учетом достаточно большого числа членов разложения гиперболических функций в~\eqref{eq:ulg:TransferFunc_TemperMody_GU_Mody_Laplas} в ряды Тейлора по степеням их аргументов. При этом, как будет показано далее, при практическом решении задачи~\eqref{eq:ulg:GU_Mody_via_TemperMody_Laplas}, как правило, переход от оригиналов к изображениям и обратно  реализуется современными программными средствами моделирования динамических систем.





\smallskip



\Section{Планирование температурных измерений}

Точность решения ОЗТ во многом определяется входной экспериментальной информацией, которая существенно зависит от расположения датчиков температуры, в связи с чем  эффективное решение двумерных обратных задач возможно только при предварительно выполненном этапе оптимального планирования температурных измерений. 

В рамках предлагаемого подхода найдем оптимальную в определенном смысле  совокупность значений $ N+1 $   координат $ y_i$, $i = \overline {1,N+1}, $  точек $\{ (x^*,y_1),(x^*,y_2), \ldots, (x^*,y_{N + 1}) \}$ контроля температуры  при заданном числе $ N  $, определяемом априори выбранным числом членов усеченного ряда~\eqref{eq:ulg:Finit_Series_GU_Laplas}.



При вариациях набора  $ (x^*,y_i)$, $i = \overline {1,N + 1}, $ точек контроля решение ${\overline \theta  _n}$, 

$n = \overline {0,N}, $  системы уравнений ~\eqref{eq:ulg:Finit_Series_Experiment} определяется в форме соответствующей зависимости $ {\overline \theta  _n}(x^*,\varphi ,\Delta ) $  от вектора  $\Delta  = ( y_i )$, 

$i = \overline {1,N + 1}, $ значений искомых координат. 

В~таком случае расчетное температурное поле представляется в~форме оригинала укороченной суммы $ N + 1 $ членов ряда~\eqref{eq:ulg:Series_Temper_Laplas}:

\begin{equation}\label{eq:ulg:Temperature_Planirovanie_Experimentov}

\theta _{N + 1}( x^*,y,\varphi ,\Delta  ) \approx \sum\limits_{n = 0}^N {C_n^2{{\overline \theta  }_n} ( x^*,\varphi ,\Delta  )}  \cos \lambda_ny,\quad y \in [ 0,1 ], \quad \varphi  \in [ 0,\varphi ^* ] .    

\end{equation}



Согласно концепции минимаксной оптимизации, реализующей оценивание температурного отклонения в равномерной метрике~\cite{ulg:bib:Rapoport}, в качестве критерия оптимизации будем использовать минимизацию максимального значения температурной невязки во всей области определения аргументов:

\begin{equation}\label{eq:ulg:Minimax_Planirovanie_Experimentov}

{I_1}( x^*,\Delta ) = 

\max \limits_{\substack{

\varphi  \in [ 0, \varphi ^*]\\

y \in [0,1]}}

 \bigl| \theta _{N + 1} (x^*,y,\varphi ,\Delta ) - 

 \theta ^* (x^*,y,\varphi ) \bigr| \to \min _{( x^*,\Delta )} ,

\end{equation}

\noindent

где  $\theta ^*(x^*,y,\varphi )$ "--- измеряемая температура, полагаемая доступной для контроля при всех $  y \in [ 0, 1]  $.



Учет очевидных соображений физического характера приводит к тому, что ошибка аппроксимации с~течением времени увеличивается и~достигает своего максимального значения в~последний момент времени  $ \varphi ^*  $, в связи с~чем критерий~\eqref{eq:ulg:Minimax_Planirovanie_Experimentov} может быть сформулирован по отношению не ко всему интервалу идентификации  $[ 0,\varphi ^* ] \ni \varphi $, а к моменту  $ \varphi ^*  $:

\begin{equation}\label{eq:ulg:Minimax_Planirovanie_Experimentov_END}

{I_2}(x^*,\Delta ) = 

\max \limits_{y \in [0,1]\,} 

\bigl| \theta _{N + 1}( x^*,y,\varphi ^*,\Delta ) - 

\theta ^*(x^*,y,\varphi ^*) \bigr| \to \min _{( x^*,\Delta  )} .    

\end{equation}



Таким образом, минимаксная задача математического программирования~\eqref{eq:ulg:Minimax_Planirovanie_Experimentov_END} может быть решена относительно координат  $ (x^*,y_i)$, $i = \overline {1,N + 1} ,$ или, при фиксированном значении  $ x ^*  $, заданном, например, исходя из  технологических требований, относительно вектора  $\Delta= (y_i)$, $i = \overline {1,N+1} $. 



Для решения задачи~\eqref{eq:ulg:Minimax_Planirovanie_Experimentov_END} могут быть использованы известные качественные свойства оптимальных температурных распределений, обладающих при оптимальных значениях координат $  \ (x^*,y_1^0),(x^*,y_2^0), \ldots, (x^*,y_{N + 1}^0)  $  наилучшим равномерным приближением к требуемому состоянию $ {\theta ^*}(x^*,y,{\varphi })$~\cite{ulg:bib:Rapoport}. 

Данные свойства фиксируют в отдельных точках $ (x^*,y_j^{al})$, $j = \overline {1,R}, $  на интервале $ y \in [ 0,1 ] $  достижение предельных отклонений расчетной температуры $ \theta _{N + 1}( x^*,y_j^{al},\varphi ^*,\Delta ^0 ) $   от экспериментальной $ \theta ^*(x^*,y_j^{al},\varphi ^*) $, равных  $ \pm {I_2}(x^*,\Delta ^0) $, и тем самым устанавливают конфигурацию погрешности  $\theta _{N + 1}( x^*,y,\varphi ^*,\Delta ^0)  \hm - \theta ^*(x^*,y,\varphi ^*)$ аппроксимации температурного распределения  $\theta ^*(x^*,y,\varphi ^*)$.



Оптимальное решение  $ \theta _{N + 1}( x^*,y,\varphi ^*,\Delta ^0) - \theta ^*(x^*,y,\varphi ^*)$, соответствующее вектору значений  $ \Delta ^0$%,\,i = \overline {1,N + 1} $ 

для каждого фиксированного  $x ^*$ , обладает свойствами чебышевского альтернанса, в соответствии с~чем при $\varphi  = \varphi ^*$  на интервале  ${y \in [ 0,1]}$ в точках  $(x^*,y_j^{al})$ , число которых $R$  на единицу превышает количество  $N + 1$  искомых координат, достигаются знакочередующиеся максимальные по величине отклонения, равные

$ \pm \bigl( \theta _{N + 1} ( x^*,y_j^{al}, \varphi ^*, \Delta ^0) \hm - \theta ^*(x^*,y_j^{al}, \varphi ^*) \bigr)$. 



Данное свойство приводит к замкнутой системе $N + 2$  соотношений для предельных значений температурных невязок в этих точках относительно неизвестных "--- значений координат  $(y_j^{al})$ точек альтернанса и величины предельного отклонения  ${I_2}( x^*,\Delta ^0 )$. 

Условия существования экстремума 

$$\frac{\partial }{\partial y}\bigl( \theta _{N + 1}( x^*,y_s^{al},\varphi ^*,\Delta ^0 ) - {\theta ^*}(x^*,y_s^{al},\varphi ^*) \bigr)$$ во внутренних точках $y_s^{al} \in \{ y_j^{al} \}$, $j = \overline {1,R}, $  интервала $[ 0,1 ] \ni y$ дополняют систему, которая в результате принимает вид

\begin{equation}\label{eq:ulg:Alternance_System}

\theta _{N + 1}( x^*,y_j^{al},\varphi ^*,\Delta ^0 ) - \theta ^*(x^*,y_j^{al},\varphi ^*) =  \pm {( - 1)^{j + 1}}{I_2};

\end{equation}

\begin{equation}\label{eq:ulg:Alternance_System_Derivative}

\frac{\partial }{\partial y}\bigl( \theta _{N + 1}( x^*,y_s^{al},\varphi ^*,\Delta ^0 ) - \theta ^*(x^*,y_s^{al},\varphi ^*) \bigr) = 0.

\end{equation}



Влияние координаты $x^*$ на точность решения задачи может быть исследовано методом математического моделирования. 



\smallskip



\Section{Пример решения обратной задачи теплопроводности}

Для демонстрации возможностей представленного метода была решена серия ОЗТ по восстановлению пространственно-временного управления $ q(y,\varphi ) $  методом модальной идентификации с учетом от трех до десяти мод температурного поля.

Модельное идентифицируемое воздействие $ q^*(y,\varphi )$ в форме произведения двух функций одной переменной $ q^*(y,\varphi ) = q^{(1)}(y)q^{(2)}(\varphi ) $  имело вид

\begin{equation}\label{eq:ulg:Seeking control}

q^*(y,\varphi ) = \sin (\alpha \pi y)\bigl( e^{\beta \varphi } - 1 \bigr),

\end{equation}

с учетом которого было получено выражение, моделирующее наблюдаемое температурное поле~\cite{ulg:bib:Koshlyakov_Gliner_Smirnov}:

\begin{multline}\label{eq:ulg:2D_Temperature}

\theta ^*(x,y,\varphi ) = \int _0^\varphi  {\overline q }_{00}(\tau )d\tau  + 

2\sum\limits_{m = 1}^\infty  \cos  {\mu _m}x  \int _0^\varphi  {{\overline q }_{m0}(\tau ){e^{ - \mu _m^2(\varphi  - \tau )}}d\tau } + 

\\

+ 2\sum\limits_{n = 1}^\infty  \cos  \lambda _n y  \int _0^\varphi  {{\overline q }_{0n}}(\tau ){e^{ - \lambda _n^2(\varphi  - \tau )}}d\tau   +

\\

+ 4\sum\limits_{m = 1}^\infty  \sum\limits_{n = 1}^\infty  

\cos \mu _m x  \cos  \lambda _n y   \int _0^\varphi  {\overline q }_{mn}(\tau ){e^{ - (\mu _m^2 + \lambda _n^2)(\varphi  - \tau )}}d\tau ,

\end{multline}



\begin{equation}\label{eq:ulg:q_MN}

{\overline q _{mn}}(\varphi ) = {\overline q _n}(\varphi )\cos \mu_m,\quad m, n = 0, 1, 2, \dots,

\end{equation}

где  $ {\overline q _{n}}(\varphi )$ определяются согласно~\eqref{eq:ulg:Mody_Control} с подстановкой  $ q(y,\varphi ) $ в форме~\eqref{eq:ulg:Seeking control}.



На основе выражений~\eqref{eq:ulg:Seeking control},~\eqref{eq:ulg:2D_Temperature} была решена серия ОЗТ для ряда возрастающих значений  $ N $.



Предварительно для каждого конкретного значения $ N $  решалась задача планирования температурных измерений~\eqref{eq:ulg:Minimax_Planirovanie_Experimentov_END} при использовании  выражения~\eqref{eq:ulg:2D_Temperature} для $ \theta ^*(x^*,y,\varphi ^*) $. При этом координата  $ x^* $ полагалась заданной, и определению подлежали координаты  $(y_i^0),\,i = \overline {1,N + 1} $, значения которых были найдены в результате численного решения системы альтернансных соотношений ~\eqref{eq:ulg:Alternance_System}.



Далее были решены прямые задачи теплопроводности, моделирующие на интервале идентификации $ \varphi  \in [ 0,\varphi ^* ] $ экспериментальные температурные зависимости $\theta ^*(x^*,{y_i},\varphi )$ в~точках с~найденными координатами $(x^*,{y_i^0})$, $i = \overline {1,N + 1} $.



Равенства~\eqref{eq:ulg:Finit_Series_Experiment} позволяют записать для всей совокупности экспериментальных кривых систему $ N $  алгебраических уравнений

\begin{equation}\label{eq:ulg:Via_Matrice_Sobstvennyh_Function}

\boldsymbol \theta ^* = \boldsymbol  \Phi   \overline {\boldsymbol \theta } ,

\end{equation}

где  $\boldsymbol \theta ^* = \bigl[ \theta ^*(x^*,y_1,\varphi )\ \ \theta ^*(x^*,{y_2},\varphi )\ \ \cdots \ \ \theta ^*(x^*,y_{N + 1},\varphi ) \bigr]^\top $ "--- вектор-столбец  температурных зависимостей, полученных во всех точках контроля на всем интервале идентификации,

матрица 

\begin{equation}\label{eq:ulg:Matrice_Sobstvennyh_Function}

\boldsymbol \Phi  = \left[ {\begin{array}{*{20}{c}}

{C_0^2}&{C_1^2  \cos \pi {y_1}}&{\cdots}&{C_{N - 1}^2 \cos (N - 1)  \pi {y_1}}\\

{C_0^2}&{C_1^2  \cos \pi {y_2}}&{\cdots}&{C_{N - 1}^2 \cos (N - 1)\pi {y_2}}\\

{\vdots}&{\vdots}&{\ddots}&{\vdots}\\

{C_0^2}&{C_1^2  \cos \pi {y_N}}&{\cdots}&{C_{N - 1}^2  \cos (N - 1) \pi {y_N}}

\end{array}} \right]

\end{equation}

составлена из значений собственных функций в точках контроля с учетом нормирующих коэффициентов, а  

$\overline {\boldsymbol \theta }  = \bigl[ 

\overline {\theta } _0 (x^*,\varphi ) \ \ \overline {\theta } _1(x^*,\varphi ) \ \ \cdots \ \ 

\overline {\theta } _{N - 1} (x^*,\varphi )

\bigr]^\top$ "--- вектор-столбец модальных составляющих температур $\theta (x^*,y_i,\varphi )$.



Таким образом, на основе~\eqref{eq:ulg:Via_Matrice_Sobstvennyh_Function} при известном векторе $\boldsymbol {\theta }^*$ могут быть определены $N+1$  модальные составляющие  $\overline \theta  _n(x^*,\varphi )$:

\begin{equation}\label{eq:ulg:Via_MNK}

\overline {\boldsymbol {\theta }}  = \boldsymbol \Phi ^{ - 1}  \boldsymbol {\theta }^*.

\end{equation}



Далее с использованием полученных значений  $ {\overline \theta  _n}(x^*,\varphi ) $ решается в пространстве изображений Лапласа задача ~\eqref{eq:ulg:GU_Mody_via_TemperMody_Laplas} восстановления модальных составляющих  ${\widetilde {\overline q }_n}(p)$ граничного управления. 

Непосредственное вычисление \linebreak ${\widetilde {\overline q }_n}(p)$~на основе~\eqref{eq:ulg:GU_Mody_via_TemperMody_Laplas}, содержащего оператор $H_n^{ - 1}(x^*,p)$, обратный непрерывному, привело бы к нарушению условия устойчивости решения, в связи с чем для реализации выражения~\eqref{eq:ulg:GU_Mody_via_TemperMody_Laplas} используется структура, представленная на рис.~\ref{Struct_Scheme}, где $ K $  "--- достаточно большое число.

Для идентификации значений всех модальных составляющих ${\widetilde {\overline q }_n}(p)$, $n \hm = \overline {0,N} $,  реализуется совокупность $ N+1 $  независимых контуров согласно структуре рис.~\ref{Struct_Scheme}.

Числитель и знаменатель передаточной функции  ${H_n}(x,p)$~\eqref{eq:ulg:TransferFunc_TemperMody_GU_Mody_Laplas}, содержащие функции гиперболического синуса и косинуса, при моделировании могут быть приближенно выражены в виде разложения в усеченные ряды по степеням~$p$:

\begin{multline}\label{eq:ulg:Hyperbolic_function}

\ch\bigl( {\sqrt {p + \pi ^2 n^2}  \cdot x} \bigr) = 1 + \frac{( {p + \pi ^2 n^2} ){x^2}}{2!} + \frac{{{( {p + \pi ^2 n^2} )}^2}{x^4}}{4!} + \cdots,

\\

\sqrt {p + \pi ^2 n^2}  \cdot \sh\sqrt {p + \pi ^2 n^2}  = ( {p + \pi ^2 n^2} ) + \frac{{( {p + {\pi ^2}{n^2}} )}^2}{3!} + \frac{{( p + \pi ^2 n^2 )}^3}{5!} + \cdots .

\end{multline}



\begin{figure}[b!]

\noindent

\begin{minipage}[h]{0.5\textwidth}



\caption {Структурная схема восстановления модальных составляющих $\widetilde {\overline q }_n (p)  $

\label{Struct_Scheme}

 }



\smallskip \footnotesize



[Figure~\ref{Struct_Scheme}.

The structure diagram of calculation \centerline{of modal components $\widetilde {\overline q }_n (p)$]}

\end{minipage} \hfill

\begin{minipage}[h]{0.5\textwidth}



\centering

 \includegraphics[width=0.96\textwidth]{Structure_cheme2}

% \includegraphics[viewport=180 370 500 470, clip=true, scale =0.9] {Structure_cheme2} \\



\noindent

% \begin{minipage}[h]{0.45\textwidth}

% \centering  \scriptsize \sl

% \includegraphics[viewport=-430 250 500 600, clip=true, scale =0.4]{Temp3D_5N}\\





\end{minipage}



\end{figure}







После нахождения составляющих ${\widetilde {\overline q }_n}(p)$  граничное воздействие $\widetilde q(y,p)$   восстанавливается в виде ряда~\eqref{eq:ulg:Series_Temper_Laplas}  с конечным числом слагаемых, оригинал которого, полученный в результате обратного преобразования Лапласа, соответствует искомой функции $q(y,\varphi )$.

При практической реализации структуры (рис.~\ref{Struct_Scheme}) в среде моделирования динамических систем {\tt Simulink MATLAB}$^\circledR$ операции прямого и обратного преобразований Лапласа производятся средствами {\tt Simulink}$^\circledR$ путем численного моделирования.





\begin{figure}[b!] 

\centering



\noindent

\begin{minipage}[h]{0.45\textwidth}

\centering  \scriptsize \sl

%\includegraphics[viewport=100 250 500 600, clip=true, scale =0.45]{Temp3D_5N}\\

\includegraphics[scale=.45]{Temp3D_5N}\\



\end{minipage}

~

\begin{minipage}[h]{0.45\textwidth}

\centering  \scriptsize \sl

%\includegraphics[viewport=100 250 500 600, clip=true, scale =0.45]{GU_3D_5N}\\

\includegraphics[scale=.45]{GU_3D_5N}\\



\end{minipage}



\caption{Температурная невязка $\theta _{N + 1}( x^*,y,\varphi ,\Delta ^0 ) - \theta ^*(x^*,y,\varphi )$, полученная в результате решения задачи планирования эксперимента (слева), и погрешность  ${q}( y,\varphi ) - q^*(y,\varphi )$ аппроксимации граничного управления (справа) при   

$ N=4 $

\label{Results_N_4} 

}



\smallskip \footnotesize



[Figure~\ref{Results_N_4}.

The temperature discrepancy $\theta _{N + 1}( x^*,y,\varphi ,\Delta ^0 ) - \theta ^*(x^*,y,\varphi )$ obtained as a result of solving the problem of experimental design (left), and the error ${q}( y,\varphi ) - q^*(y,\varphi )$

\centerline {  of approximation of the boundary control (right) for  $ N=4 $   ]}

%\end{figure}



\bigskip



%

%

%\begin{figure}[h!] 

\centering



\centering  \scriptsize 

\includegraphics[scale =0.54]{TTT-Control}\\





\caption{Температурная невязка $\theta _{N + 1}(x^*,y,\varphi ^*) - \theta ^*(x^*,y,\varphi ^*)$, полученная в результате решения задачи планирования эксперимента (слева), и погрешность $q( y,\varphi ^* ) - q^*(y,\varphi ^*)$ \centerline{аппроксимации граничного управления (справа): 

{\sl 1} "--- при  $N=2$; 

{\sl 2} "--- при $N=4$; }

\mbox{{\sl 3} "---  при  $N=7$}

\label{Discrepancy_N_2_4_7} 

}



\smallskip \footnotesize



[Figure~\ref{Discrepancy_N_2_4_7}.

The temperature discrepancy $\theta _{N + 1}(x^*,y,\varphi ^*) - \theta ^*(x^*,y,\varphi ^*)$ obtained as a result of solving the problem of experimental design (left) and the error $q( y,\varphi ^* ) - q^*(y,\varphi ^*)$ of  \centerline{approximation of the boundary control (right) 

 ({\sl 1}---for  $N=2$; {\sl 2}---for $N=4$; {\sl 3}---for  $N=7$)]}

\end{figure}





Некоторые результаты решения ОЗТ методом модальной идентификации представлены на рис.~\ref{Results_N_4},~\ref{Discrepancy_N_2_4_7}.

В таблице  приведены погрешность  

$$\varepsilon _\theta  =  \max \limits_y | \theta _{N + 1}(x^*,y,\varphi ^*) - \theta ^*(x^*,y,\varphi ^*) |$$ приближения температуры 

и ошибка $${\varepsilon _u} = \max \limits_y | q ( y,\varphi ^* ) - q^*(y,\varphi ^*) |$$  аппроксимации граничного управления для различных $N$. 

Результаты, отраженные на рисунках и в таблице, соответствуют данным $\varphi ^* = 2$, $x^* = 0.96$ и~следующим оптимальным значениям $y_i^0$ вектора $y^0=(y_i^0)$, $i = \overline {1,N + 1}$,  координат $y_i$ точек контролируемых температур: \pagebreak





$y^0 = [0.11;\,0.42;\,0.80]$ при $N = 2$; 



$y^0 = [0.08;\,0.30;\,0.57;\,0.85]$ при $N = 3$;



$y^0 = [0.06;\,0.23;\,0.43;\,0.66;\,0.88]$ при $N = 4$; 



$y^0 = [0.05;\,0.19;\,0.35;\,0.53;\,0.71;\,0.91]$ при $N = 5$; 



$y^0 = [0.04;\,0.14;\,0.26;\,0.39;\,0.52;\,0.66;\,0.80;\,0.93]$ при $N = 7$;



$y^0 = [0.03;\,0.11;\,0.20;\,0.30;\,\,0.40;\,0.51;\,0.62;\,0.73;\,0.84;\,0.95]$ при $N = 9$.





%-------------------------------------------------------



\begin{table}[t!]

\small \centering



\renewcommand{\tabcolsep}{0.55cm}

\begin{tabular}{c||c|c|c|c|c|c}

\hline

%\rule{0mm}{11pt}%

\footnotesize $N$  & \footnotesize  $2$ & \footnotesize  $3$ & \footnotesize  $4$ & \footnotesize  $5$ & \footnotesize  $7$ & \footnotesize  $9 $ \\

            \hline

            \rule{0mm}{12pt}%

            \phantom{1} ${\varepsilon _\theta }, \%$ & 0.41 & 0.16 & 0.08 & 0.04 & 0.02  & 0.01 \\

             \phantom{1} ${\varepsilon _u}, \%$ & 11.6 & 8.4 & 7.0 & 6.4 & 5.8 & 5.6 \\

             \hline

        \end{tabular}

\footnotesize 



The error $\varepsilon _\theta  =  \max \limits_y | \theta _{N + 1}(x^*,y,\varphi ^*) - \theta ^*(x^*,y,\varphi ^*) |$ of the temperature approximation  and the error   ${\varepsilon _u} = \max \limits_y | q ( y,\varphi ^* ) - q^*(y,\varphi ^*) |$ of approximation of the boundary control for different $N$ 



\end{table} 





\smallskip



%-------------------------------------------------------





\Section[N]{Заключение}

Предложенный в работе метод решения обратной задачи теплопроводности, обеспечивающий восстановление пространственно"=временного граничного управляющего воздействия, позволяет найти искомую характеристику на основе аналитической идентификации его модальных составляющих. 

Погрешность решения задачи определяется числом учитываемых слагаемых усеченного ряда и уменьшается с его увеличением.

Предварительное решение задачи оптимального планирования эксперимента позволяет определить координаты точек контроля температуры, обеспечивающие решение ОЗТ с удовлетворительной точностью.



% Не забудьте сгенерировать страницу с информацией о статье на английском языке с помощью команды \verb"\makeentitle".



\bigskip





% \AdditionalContent{Конкурирующие интересы}{Укажите какие конфликты интересов имеются.} 

\AdditionalContent{Конкурирующие интересы}{Конкурирующих интересов не имеем.} 

% \AdditionalContent{Авторский вклад и ответственность}{Укажите степень авторской ответственности каждого автора.} 

%\AdditionalContent{Авторская ответственность}{Я несу полную ответственность за предоставление окончательной версии рукописи в печать. Окончательная версия рукописи мною одобрена.} 

\AdditionalContent{Авторский вклад и ответственность}{Все  авторы  принимали  участие  в  разработке  концепции  \mbox{статьи}  и  в  написании  рукописи. Авторы несут  полную  ответственность  за  предоставление  окончательной рукописи в печать. Окончательная  версия рукописи была одобрена всеми авторами.} 

% \AdditionalContent{Финансирование}{Укажите источник финансирования работы.}

% \AdditionalContent{Финансирование}{Исследование выполнялось без финансирования.}

 \AdditionalContent{Финансирование}{Работа выполнена при поддержке Российского фонда фундаментальных исследований (проекты 18--08--00048\_a и 18--08--00565\_a).}

% \AdditionalContent{Благодарность}{Выразите свою благодарность.}

%\AdditionalContent{Благодарность}{Авторы благодарны рецензенту за тщательное прочтение статьи и~ценные предложения и комментарии.}



\begin{thebibliography}{99}



\Bibitem{ulg:bib:Ozisik}

\by \"Ozisik~M.~N., Orlande~H.~R.~B.

\book Inverse Heat Transfer: fundamentals and applications

\yr 2000

\publ Taylor and Francis

\publaddr New York

\totalpages xviii+330~p \nofrills

\crossref{10.1201/9780203749784}



\RBibitem{ulg:bib:Matsevityi_Param_funct_identif}

\by Мацевитый~Ю.~М., Гайшун~И.~В., Борухов~В.~Т., Костиков~А.~О.

\paper Параметрическая и функциональная идентификация тепловых процессов

\jour Пробл. машиностроения

\yr 2011

\vol 14

\issue 3

\pages 40--47

%\elink \url{http://journals.uran.ua/jme/article/view/74714}

%\misctransl

%\by Matsevityi~Y.~M. , Gaishun~I.~V., Borukhov~V.~T., Kostikov~A.~O.

%\paper Parametric and functional identification of thermal processes

%\jour Problems of mechanical engineering

%\yr 2011

%\vol 14

%\issue 3

%\pages 40--47

%\lang In Russian





\Bibitem{ulg:bib:Alifanov_Opening_Lecture}

\by Alifanov~O.

\paper Inverse problems in identification and modeling of thermal processes: Theory and practice (Invited paper, Opening keynote lecture)

\inbook Book of Abstracts

\bookinfo 8th International Conference on Inverse Problems in Engineering (ICIPE 2014)

\yr 2014

\pages 3--4

\publ Institute of Thermal Technology Silesian University of Technology

\publaddr Gliwice--Krakow, Poland









\RBibitem{ulg:bib:Alifanov_OZT}

\by Алифанов~О.~М.

\book Обратные задачи теплообмена

\yr 1988

\publ Машиностроение

\publaddr М.

\totalpages 280

%\transl

%\by Alifanov~Oleg~M.

%\book Inverse Heat Transfer Problems

%\yr 1994

%\publ Springer

%\publaddr Berlin, Heidelberg

%\totalpages 348

%\crossref{10.1007/978-3-642-76436-3}



\Bibitem{ulg:bib:Beck_Blackwel}

\by Beck~J.~V., Blackwell~B., St. Clair~C.~R., jr.

\book Inverse Heat Conduction: Ill-Posed Problems

\yr 1985

\publ J. Wiley and Sons

\publaddr New York

\totalpages xvii+308





\RBibitem{ulg:bib:Vatulyan_IP_MDTT}

\by Ватульян~А.~О.

\book Обратные задачи в механике деформируемого твердого тела

\yr 2007

\publ Физматлит

\publaddr М.

\totalpages 224

%\misctransl

%\by Vatul'yan~A.~O.

%\book Obratnye zadachi v mekhanike deformiruemogo tverdogo tela \rm [Inverse Problems in the Mechanics of Deformable Solids]

%\yr 2007

%\publ Fizmatlit

%\publaddr Moscow

%\lang In Russian

%\totalpages 224





\Bibitem{ulg:bib:Alifanov_Russian_contribution}

\by Alifanov~O.

\paper Inverse problems in identification and modeling of thermal processes: Russian contributions

\jour  Int. J. Numer. Methods Heat Fluid Flow

\vol 27 

\issue  3

\pages 711--728

\crossref{10.1108/HFF-03-2016-0099}



\RBibitem{ulg:bib:Matsevityi_nonlinear_IHCP}

\by Мацевитый~Ю.~М., Костиков~А.~О., Сафонов~Н.~А., Ганчин~В.~В.

\paper К решению нестационарных нелинейных граничных обратных задач теплопроводности

\jour Пробл. машиностроения

\yr 2017

\vol 20

\issue 4

\pages 15--23

%\elink \url{http://journals.uran.ua/jme/article/view/120553}

%\transl

%\by Matsevityi~Y.~M., Kostikov~A.~O., Safonov~N.~A., Ganchin~V.~V.

%\paper To the solution of non-stationary nonlinear boundary inverse heat conduction  problems

%\jour Problems of mechanical engineering

%\yr 2017

%\vol 20

%\issue 4

%\pages 15--23

%\lang In Russian



\Bibitem{ulg:bib:Zhi_Qian_Numerical_solution}

\by Zhi~Qian, Xiaoli~Feng

\paper Numerical solution of a 2D inverse heat conduction problem

\jour Inverse Problems in Science and Engineering

\yr 2013

\vol 21

\issue 3

\pages 467--484

\crossref{10.1080/17415977.2012.712526}



\Bibitem{ulg:bib:Ching_yu_boundary_estimation}

\by Ching-yu~Yang, Cha'o-Kuang~Chen

\paper The boundary estimation in two-dimensional inverse heat conduction problems

\jour J. Phys. D. Appl. Phys.

\yr 1996

\vol 29

\issue 2

\pages 333

\crossref{10.1088/0022-3727/29/2/009}



\RBibitem{ulg:bib:Kuzin_2_D}

\by Кузин~А.~Я.

\paper Регуляризованное численное решение нелинейной двумерной обратной задачи теплопроводности 

\jour ПМТФ

\yr 1995

\issue 1

\pages 106--112

%\elink \url{http://www.sibran.ru/journals/issue.php?ID=119958&ARTICLE_ID=133040}

%\transl

%\by Kuzin~A.~Ya.

%\paper Regularized numerical solution of the nonlinear, two-dimensional, inverse heat-conduction problem

%\jour Journal of Applied Mechanics and Technical Physics

%\yr 1995

%\vol 36

%\issue 1

%\pages 98--104

%\crossref{10.1007/BF02369679}



\RBibitem{ulg:bib:Tikhonov_Arsenin}

\by Тихонов~А.~Н., Арсенин~В.~Я. 

\book Методы решения некорректных задач

\yr 1979

\publ Наука

\publaddr М.

\totalpages 285

%\misctransl

%\by Tikhonov~A.~N., Arsenin~V.~Ya.

%\book Metody resheniya nekorrektnykh zadach \rm [Methods of solving ill-posed problems]

%\yr 1979

%\publ Nauka

%\publaddr Moscow

%\lang In Russian

%\totalpages 285



\RBibitem{ulg:bib:Ivanov}

\by Иванов~В.~К., Васин~В.~В., Танана~В.~П. 

\book Теория линейных некорректных задач и ее приложения

\yr 1978

\publ Наука

\publaddr М.

\totalpages 206

%\transl

%\by Ivanov~V.~K., Vasin~V.~V., Tanana~V.~P.

%\book Theory of Linear Ill-posed Problems and Its Applications

%\yr 2002

%\publ VSP

%\publaddr Utrecht

%\totalpages xiii+281 p.



\RBibitem{ulg:bib:Tikhonov}

\by Тихонов~А.~Н., Гончарский~А.~В., Степанов~В.~В., Ягола~А.~Г.

\book Регуляризирующие алгоритмы и априорная информация

\yr 1983

\publ Наука

\publaddr М.

\totalpages 200

%\misctransl

%\by Tikhonov~A.~N., Goncharsky~A.~V., Stepanov~V.~V., Yagola~A.~G. 

%\book Regulyariziruyushchie algoritmy i apriornaya informatsiya \rm [Regularizing algorithms and a priori information]

%\yr 1983

%\publ Nauka

%\publaddr Moscow

%\lang In Russian

%\totalpages 200



\RBibitem{ulg:bib:Butkovskiy}

\by Бутковский~А.~Г. 

\book Методы управления системами с распределенными параметрами

\yr 1975

\publ Наука

\publaddr М.

\totalpages 568

%\misctransl

%\by Butkovskiy A.G.

%\book Metody upravleniya sistemami s raspredelennymi parametrami \rm [Control methods for distributed parameter systems]

%\yr 1975

%\publ Nauka

%\publaddr Moscow

%\lang In Russian

%\totalpages 568



\Bibitem{ulg:bib:Reinhardt}

\by Reinhardt~H.-J.

\paper A numerical method for the solution of two-dimensional inverse heat conduction problems

\jour Int. J. Numer. Methods Eng.

\yr 1991

\vol 32

\issue 2

\pages  363--383

\crossref{10.1002/nme.1620320209}



\RBibitem{KolForKuz15}

\by Колесник~С.~А., Формалев~В.~Ф., Кузнецова~Е.~Л.

\paper О граничной обратной задаче теплопроводности по восстановлению тепловых потоков к границам анизотропных тел

\jour ТВТ

\yr 2015

\vol 53

\issue 1

\pages 72--77

\mathnet{http://mi.mathnet.ru/tvt3914}

\crossref{10.7868/S0040364415010111}

\elib{http://elibrary.ru/item.asp?id=22841002}

%\transl

%\jour High Temperature

%\yr 2015

%\vol 53

%\issue 1

%\pages 68--72

%\crossref{10.1134/S0018151X15010113}

%\isi{http://gateway.isiknowledge.com/gateway/Gateway.cgi?GWVersion=2&SrcApp=PARTNER_APP&SrcAuth=LinksAMR&DestLinkType=FullRecord&DestApp=ALL_WOS&KeyUT=000350357800011}

%\elib{http://elibrary.ru/item.asp?id=24010809}

%\scopus{http://www.scopus.com/record/display.url?origin=inward&eid=2-s2.0-84923830469}





\RBibitem{ulg:bib:Koshlyakov_Gliner_Smirnov}

\by Кошляков~Н.~С., Глинер~Э.~Б., Смирнов~М.~М. 

\book Уравнения в частных производных математической физики

\yr 1970

\publ Высш. шк.

\publaddr М.

\totalpages 712

%\misctransl

%\by Koshlyakov~N.~S., Gliner~E.~B., Smirnov~M.~M. 

%\book Uravneniya v chastnykh proizvodnykh matematicheskoy fiziki \rm [The partial derivatives equations in mathematical physics]

%\yr 1970

%\publ Vysshaya Shkola

%\publaddr Moscow

%\lang In Russian

%\totalpages 712





\RBibitem{ulg:bib:Rapoport_Open-loop_controllability}

\by Рапопорт~Э.~Я.

\paper Программная управляемость линейных многомерных систем с распределенными параметрами

\jour Изв. РАН. Теория и системы управления

\yr 2015

\issue 2

\pages 22--39

\crossref{10.7868/S0002338815020110}

%\transl

%\by Rapoport~E.~Ya.

%\paper Open-loop controllability of multidimensional linear systems with distributed parameters

%\jour Journal of Computer and Systems Sciences International

%\jour J. Comput. Syst. Sci. Int.

%\yr 2015

%\vol 54

%\issue 2

%\pages 184--201

%\crossref{10.1134/S1064230715020112}



\RBibitem{ulg:bib:Rapoport}

\by Рапопорт~Э.~Я. 

\book Альтернансный метод в прикладных задачах оптимизации

\yr 2000

\publ Наука

\publaddr М.

\totalpages 336

%\misctransl

%\by Rapoport~E.~Ya. 

%\book Al'ternansnyy metod v prikladnykh zadachakh optimizatsii \rm [Alternance method in applied optimization problems]

%\yr 2000

%\publ Nauka

%\publaddr Moscow

%\lang In Russian

%\totalpages 336





\end{thebibliography}





\newpage



\makeentitle



\vspace{-5mm}



\AdditionalContent{Competing interests}{We declare that we have no conflicts of interests with the authorship and publication of this article.} 

\AdditionalContent{Authors' contributions and responsibilities}{Each author has participated in the article concept development and in the manuscript writing. The authors are absolutely responsible for submitting the final manuscript in print. Each author has approved the final version of manuscript.} 



\AdditionalContent{Funding}{This work was supported by the Russian Foundation for Basic Research (projects nos. 18--08--00048\_a and 18--08--00565\_a).}







\begin{thebibliography}{99}



\Bibitem{ulg:bib:Ozisik:ew}

\by \"Ozisik~M.~N., Orlande~H.~R.~B.

\book Inverse Heat Transfer: fundamentals and applications

\yr 2000

\publ Taylor and Francis

\publaddr New York

\totalpages xviii+330~p \nofrills

\crossref{10.1201/9780203749784}



\Bibitem{ulg:bib:Matsevityi_Param_funct_identif:en}

\by Matsevityi~Y.~M. , Gaishun~I.~V., Borukhov~V.~T., Kostikov~A.~O.

\paper Parametric and functional identification of thermal processes

\jour Journal of Mechanical Engineering

\yr 2011

\vol 14

\issue 3

\pages 40--47

\lang In Russian





\Bibitem{ulg:bib:Alifanov_Opening_Lecture:en}

\by Alifanov~O.

\paper Inverse problems in identification and modeling of thermal processes: Theory and practice (Invited paper, Opening keynote lecture)

\inbook Book of Abstracts

\bookinfo 8th International Conference on Inverse Problems in Engineering (ICIPE 2014)

\yr 2014

\pages 3--4

\publ Institute of Thermal Technology Silesian University of Technology

\publaddr Gliwice--Krakow, Poland









\Bibitem{ulg:bib:Alifanov_OZT:en}

\by Alifanov~O.~M.

\book Inverse Heat Transfer Problems

\yr 1994

\publ Springer

\publaddr Berlin, Heidelberg

\totalpages xii+348~p \nofrills

\crossref{10.1007/978-3-642-76436-3}



\Bibitem{ulg:bib:Beck_Blackwel:en}

\by Beck~J.~V., Blackwell~B., St. Clair~C.~R., jr.

\book Inverse Heat Conduction: Ill-Posed Problems

\yr 1985

\publ J. Wiley and Sons

\publaddr New York

\totalpages xvii+308





\Bibitem{ulg:bib:Vatulyan_IP_MDTT:en}

\by Vatul'yan~A.~O.

\book Obratnye zadachi v mekhanike deformiruemogo tverdogo tela \rm [Inverse Problems in  Mechanics of Deformable Solids]

\yr 2007

\publ Fizmatlit

\publaddr Moscow

\lang In Russian

\totalpages 224





\Bibitem{ulg:bib:Alifanov_Russian_contribution:en}

\by Alifanov~O.

\paper Inverse problems in identification and modeling of thermal processes: Russian contributions

\jour  Int. J. Numer. Methods Heat Fluid Flow

\vol 27 

\issue  3

\pages 711--728

\crossref{10.1108/HFF-03-2016-0099}



\Bibitem{ulg:bib:Matsevityi_nonlinear_IHCP:en}

\by Matsevityi~Y.~M., Kostikov~A.~O., Safonov~N.~A., Ganchin~V.~V.

\paper To the solution of non-stationary nonlinear reverse problems of thermal conductivity

\jour Journal of mechanical engineering

\yr 2017

\vol 20

\issue 4

\pages 15--23

\lang In Russian



\Bibitem{ulg:bib:Zhi_Qian_Numerical_solution:en}

\by Zhi~Qian, Xiaoli~Feng

\paper Numerical solution of a 2D inverse heat conduction problem

\jour Inverse Problems in Science and Engineering

\yr 2013

\vol 21

\issue 3

\pages 467--484

\crossref{10.1080/17415977.2012.712526}



\Bibitem{ulg:bib:Ching_yu_boundary_estimation:en}

\by Ching-yu~Yang, Cha'o-Kuang~Chen

\paper The boundary estimation in two-dimensional inverse heat conduction problems

\jour J. Phys. D. Appl. Phys.

\yr 1996

\vol 29

\issue 2

\pages 333

\crossref{10.1088/0022-3727/29/2/009}



\Bibitem{ulg:bib:Kuzin_2_D:en}

\by Kuzin~A.~Ya.

\paper Regularized numerical solution of the nonlinear, two-dimensional, inverse heat-conduction problem

\jour J. Appl. Mech. Tech. Phys.

\yr 1995

\vol 36

\issue 1

\pages 98--104

\crossref{10.1007/BF02369679}



\Bibitem{ulg:bib:Tikhonov_Arsenin:en}

\by Tikhonov~A.~N., Arsenin~V.~Ya.

\book Solutions of ill-posed problems

\zmath{0354.65028}

\serial Scripta Series in Mathematics

\publaddr New York 

\publ John Wiley \& Sons

\totalpages xiii+258 

\yr 1977



\Bibitem{ulg:bib:Ivanov:en}

\by Ivanov~V.~K., Vasin~V.~V., Tanana~V.~P.

\book Theory of linear ill-posed problems and its applications

\zmath{1037.65056}

\serial Inverse and Ill-Posed Problems Series

\publaddr Utrecht

\publ VSP 

\totalpages xiii+281 

\yr 2002

\vol 36









\Bibitem{ulg:bib:Tikhonov:en}

\by Tikhonov~A.~N., Goncharsky~A.~V., Stepanov~V.~V., Yagola~A.~G. 

\book Reguliariziruiushchie algoritmy i apriornaia informatsiia \rm [Regularizing algorithms and a priori information]

\yr 1983

\publ Nauka

\publaddr Moscow

\lang In Russian

\totalpages 200



\Bibitem{ulg:bib:Butkovskiy:en}

\by Butkovskiy~A.~G.

\book Metody upravleniia sistemami s raspredelennymi parametrami \rm [Control methods for systems with distributed parameters]

\yr 1975

\publ Nauka

\publaddr Moscow

\lang In Russian

\totalpages 568



\Bibitem{ulg:bib:Reinhardt}

\by Reinhardt~H.-J.

\paper A numerical method for the solution of two-dimensional inverse heat conduction problems

\jour Int. J. Numer. Methods Eng.

\yr 1991

\vol 32

\issue 2

\pages  363--383

\crossref{10.1002/nme.1620320209}



\Bibitem{KolForKuz15:en}

\paper On inverse boundary thermal conductivity problem of recovery of heat fluxes to the boundaries of anisotropic bodies

\by Kolesnik~S.~A.,  Formalev~V.~F.,  Kuznetsova~E.~L.

\jour High Temperature

\yr 2015

\vol 53

\issue 1

\pages 68--72

\crossref{10.1134/S0018151X15010113}

\isi{http://gateway.isiknowledge.com/gateway/Gateway.cgi?GWVersion=2&SrcApp=PARTNER_APP&SrcAuth=LinksAMR&DestLinkType=FullRecord&DestApp=ALL_WOS&KeyUT=000350357800011}

\elib{http://elibrary.ru/item.asp?id=24010809}

\scopus{http://www.scopus.com/record/display.url?origin=inward&eid=2-s2.0-84923830469}





\Bibitem{ulg:bib:Koshlyakov_Gliner_Smirnov:en}

\by Koshlyakov~N.~S., Gliner~E.~B., Smirnov~M.~M. 

\book Differential Equations of Mathematical Physics

\zmath{0115.30701}

\publaddr  Amsterdam

\publ North-Holland Publ.

\totalpages xvi+701 

\yr 1964





\Bibitem{ulg:bib:Rapoport_Open-loop_controllability:en}

\by Rapoport~E.~Ya.

\paper Open-loop controllability of multidimensional linear systems with distributed parameters

\jour J. Comput. Syst. Sci. Int.

\yr 2015

\vol 54

\issue 2

\pages 184--201

\crossref{10.1134/S1064230715020112}



\Bibitem{ulg:bib:Rapoport:e}

\by Rapoport~E.~Ya. 

\book Al'ternansnyy metod v prikladnykh zadachakh optimizatsii \rm [Alternance method in applied optimization problems]

\yr 2000

\publ Nauka

\publaddr Moscow

\lang In Russian

\totalpages 336





\end{thebibliography}





%\end{document}

